\clearpage
\section{Chronological Postconciliar Change Table (1963--2021)}
Principal promulgated acts governing postconciliar sanctuary furnishings, sacred vessels, linens, books, vesture, and ministry distinctions from the Second Vatican Council through the 2011 United States implementation and its 2018--2021 official companions and emendations.

{\footnotesize\begin{longtable}{>{\raggedright\arraybackslash}p{0.15\linewidth}>{\raggedright\arraybackslash}p{0.33\linewidth}>{\raggedright\arraybackslash}p{0.45\linewidth}}
\toprule
\textbf{Date} & \textbf{Promulgated act or edition} & \textbf{Sanctuary and object effect} \\
\midrule
\endhead
\textbf{4 Dec 1963} & Sacrosanctum Concilium (Second Vatican Council) & Mandates revision of the Roman Missal (no. 50), richer proclamation of Scripture (nos. 35, 51), noble simplicity in rites, vestments, and art (nos. 34, 122--128), and extension of Communion under both kinds (no. 55). \\
\textbf{26 Sept 1964} & Inter Oecumenici (SCR / Consilium; effective 7 Mar 1965) & First structural implementation of conciliar reform; recommends freestanding altar built separate from the wall (no. 91) and introduces ambo, chair, and tabernacle adjustments. \\
\textbf{4 May 1967} & Tres abhinc annos (SCR / Consilium; effective 29 June 1967) & Declares the maniple no longer required (no. 25); reduces altar kisses, genuflections, and signs of the cross. \\
\textbf{3 Apr 1969} & Missale Romanum (Ap. Const., Paul VI; Order of Mass effective 30 Nov 1969) & Promulgates the revised Roman Missal, three additional Eucharistic Prayers, and the three-year Sunday / two-year weekday Lectionary separation. \\
\textbf{26 Mar 1970} & Missale Romanum, editio typica prima (Decree Celebrationis eucharisticae, Prot. 166/70) & First typical edition of postconciliar Missal and IGMR. Briefly retains transitional subdeacon (nos. 142--152) and tunicle (no. 301); omits chalice veil mention and explicit corpus on altar cross; uses broad vessel-material wording. \\
\textbf{3 Dec 1970} & Ordo benedicendi oleum (SCDW; editio typica 1971) & Revises Chrism Mass and blessing of Holy Oils; permits vegetable oils in addition to olive oil. \\
\textbf{15 Aug 1972} & Ministeria quaedam (Motu Proprio, Paul VI; effective 1 Jan 1973) & Suppresses first tonsure, minor orders, and the major order of subdeacon in the Latin Church; establishes the two lay instituted ministries of Lector and Acolyte. \\
\textbf{29 Jan 1973} & Immensae caritatis (SCDS) & Establishes universal discipline for Extraordinary Ministers of Holy Communion in genuine pastoral need. \\
\textbf{21 June 1973} & De sacra communione et de cultu mysterii eucharistici extra Missam (SCDW) & Publishes postconciliar norms for Communion outside Mass, Eucharistic exposition, processions, and Benediction (retaining white cope and humeral veil). \\
\textbf{27 Mar 1975} & Missale Romanum, editio typica altera (Decree Cum Missale Romanum) & Removes subdeacon and tunicle from IGMR, substituting instituted acolyte and lector; retains 1975 nos. 270, 276, and 290--291 wording on cross, tabernacle, and vessel materials. \\
\textbf{21 Jan 1981} & Ordo lectionum Missae, editio typica altera (SCSD) & Expanded General Introduction to the Lectionary affirming dignity of the Lectionary and strongly commending a separate Book of the Gospels (nos. 35--37). \\
\textbf{25 Jan 1983} & Codex Iuris Canonici (John Paul II; effective 27 Nov 1983) & Codifies universal law on lay ministries (can. 230), holy oils (can. 847), parish baptismal font (can. 858), ministers of Communion (can. 910), tabernacle, pyx, and sanctuary lamp (cann. 938--940), and fixed/movable altars (cann. 1235--1239). \\
\textbf{14 Sept 1984} & Caeremoniale Episcoporum, editio typica (CDW) & Codifies episcopal and solemn ceremonial appointments, Chrism Mass (no. 274), and vestment norms including that the surplice is always worn over a cassock (no. 65). \\
\textbf{20 Apr 2000 / 2002} & Missale Romanum, editio typica tertia \& CDWDS Response (25 Sept 2000, Prot. N. 2036/00/L) & Restores mandatory figure of Christ crucified on the cross (GIRM 308), precious metal as primary vessel norm (328--330), explicit praise of the optional chalice veil (118), two co-equal tabernacle options judged by the bishop (314--315), and clarifies GIRM 299 freestanding altar wording. \\
\textbf{22 Mar 2002} & USCCB Norms for Holy Communion under Both Kinds (Prot. 1383/01/L; revised 2011) & Establishes particular law for U.S. dioceses on Communion under both kinds; paired initially with a 3-year EMHC purification indult whose non-renewal CDWDS confirmed in October 2006. \\
\textbf{25 Mar 2004} & Redemptionis Sacramentum (CDWDS; AAS 96 [2004] 549--601) & Strictly prohibits post-consecration pouring and flagons for the Precious Blood (no. 106), reprobates glass, earthenware, and clay sacred vessels (no. 117), mandates sacrarium first-rinsing of linens (no. 120), and guards ordained vesture and EMHC limits (nos. 121--128, 154--160). \\
\textbf{2008} & Missale Romanum, editio typica tertia, reimpressio emendata (Notitiae 44 [2008] 367--388) & Emended Latin reprint incorporating textual corrections, new dismissals, and Supplementum; IGMR sanctuary, vessel, linen, book, and vesture provisions remain unchanged. \\
\textbf{27 Nov 2011} & The Roman Missal, Third Edition (United States of America) & Implements third typical edition in U.S. dioceses with CDWDS-confirmed U.S. GIRM adaptations (nos. 283, 301, 304, 329, 339, 346e, 346h) and harmonized 2011 USCCB Communion Norms. \\
\textbf{2018 / 2021} & USCCB Excerpts from the Roman Missal (2018), Spiritus Domini (10 Jan 2021), and Postquam Summus Pontifex (22 Oct 2021) & Publishes official U.S. Book of the Chair (2018); amends CIC can. 230 §1 to open instituted Lector and Acolyte to lay men and women (Spiritus Domini); adjusts liturgical-book preliminary norms to amended can. 838 (Prot. n. 394/21, incorporated with Prot. n. 228/20 in the U.S. online GIRM). \\
\bottomrule
\end{longtable}}

\clearpage
\section{Keyed Source Notes}
Each bracketed key on the plates resolves to the governing edition loci and 1962 comparison witnesses below.

{\footnotesize\begin{longtable}{>{\raggedright\arraybackslash}p{0.08\linewidth}>{\raggedright\arraybackslash}p{0.24\linewidth}>{\raggedright\arraybackslash}p{0.61\linewidth}}
\toprule
\textbf{Key} & \textbf{Entry} & \textbf{Governing loci and 1962 comparison} \\
\midrule
\endhead
\textbf{N01} & \textbf{Altar} & GIRM 273--274, 280, 296--308; USCCB GIRM 301, 304; CIC (1983) cann. 1235--1239; CDWDS Response Prot. N. 2036/00/L (25 Sept 2000). 1962 counterpart: Missale Romanum (1962), Rubricae Generales 526--527; Ritus Servandus II.2, IV.5, V.3. Symbolism: GIRM 296, 298; 1 Pet 2:4; 1 Cor 10:4; Heb 13:10; St. Thomas Aquinas, Summa Theologiae III, q. 83, a. 3. \\
\textbf{N02} & \textbf{Ambo} & GIRM 28, 118b, 130--136, 175, 196, 309; Ordo lectionum Missae (1981) Praenotanda 32--34; Inter Oecumenici (1964) 96. 1962 counterpart: Missale Romanum (1962), Ritus Servandus VI.4--5 (no fixed sanctuary ambo prescribed). Symbolism: GIRM 28, 309; Ordo lectionum Missae (1981) Praenotanda 32--34; Ps 118:105 [Heb 119:105]. \\
\textbf{N03} & \textbf{Presidential Chair} & GIRM 118a, 124, 127--130, 165--167, 310; Inter Oecumenici (1964) 92; Caeremoniale Episcoporum (1984). 1962 counterpart: Missale Romanum (1962), Ritus Servandus VI.4; Caeremoniale Episcoporum (1948) I.12--13; Fortescue (1943), pp. 8--9. Symbolism: GIRM 310. \\
\textbf{N04} & \textbf{Credence Table} & GIRM 118c, 139, 163, 178, 190, 192, 279, 306; Redemptionis Sacramentum (2004) 119; USCCB Norms for Communion under Both Kinds (2011) 53. 1962 counterpart: Missale Romanum (1962), Rubricae Generales 528; Ritus Servandus II.5; Caeremoniale Episcoporum (1948) I.12. \\
\textbf{N05} & \textbf{Tabernacle} & GIRM 274, 314--317; CIC (1983) can. 938; Redemptionis Sacramentum (2004) 130--131; Eucharisticum Mysterium (1967) 52--55. 1962 counterpart: Codex Iuris Canonici (1917) cann. 1268--1269; Missale Romanum (1962), Ritus Servandus II.2. Symbolism: GIRM 314--315; Eucharisticum Mysterium (1967) 53--55; Ex 25:8; St. Thomas Aquinas, Summa Theologiae III, q. 83, a. 5, ad 11. \\
\textbf{N06} & \textbf{Sanctuary Lamp} & GIRM 316; CIC (1983) can. 940. 1962 counterpart: Codex Iuris Canonici (1917) can. 1271. Symbolism: GIRM 316; CIC (1983) can. 940; Ex 27:20--21; Lev 24:2--3. \\
\textbf{N07} & \textbf{Baptismal Font} & CIC (1983) can. 858, can. 861 §1; Ordo Baptismi Parvulorum (General Introduction 25, 45--47); Roman Missal (2008/2011) Easter Vigil Blessing of Baptismal Water; CCC 1237--1243. 1962 counterpart: Missale Romanum (1962), Vigilia Paschalis (cf. Missale Romanum 1922, Holy Saturday); Rituale Romanum (1925/1952), Tit. II. Symbolism: Roman Missal Easter Vigil; CCC 1237--1243; John 3:5; Rom 6:3--4; St. Cyril of Jerusalem, Mystagogical Catechesis 2.4. \\
\textbf{N08} & \textbf{Chalice} & GIRM 118c, 139, 142, 150--151, 163, 279, 306, 327--330, 333; Redemptionis Sacramentum 106, 117--119; USCCB Norms for Holy Communion under Both Kinds (2011) 35, 52--55. 1962 counterpart: Missale Romanum (1962), Ritus Servandus I.1, II.1, VIII. Symbolism: GIRM 327; Matt 26:27--28; 1 Cor 10:16; Ps 115:13 [Heb 116:13]; St. Thomas Aquinas, Summa Theologiae III, q. 83, a. 3, ad 6 \& a. 5, ad 9. \\
\textbf{N09} & \textbf{Paten} & GIRM 118c, 139, 155, 278--279, 306, 327--329, 331, 333; USCCB GIRM 329; Redemptionis Sacramentum 93, 117--119. 1962 counterpart: Missale Romanum (1962), Ritus Servandus I.1, VII.9, X.1--2, X.7. Symbolism: GIRM 327, 331; 1 Cor 10:17; St. Thomas Aquinas, Summa Theologiae III, q. 83, a. 3, obj. 6 \& ad 6. \\
\textbf{N10} & \textbf{Ciborium} & GIRM 118c, 163, 279, 306, 327--329, 333; USCCB GIRM 329; CIC (1983) can. 939; Redemptionis Sacramentum 117--119; USCCB Norms for Holy Communion under Both Kinds (2011) 37, 53. 1962 counterpart: Missale Romanum (1962), Ritus Servandus II.3, VII.3, X.6; CIC (1917) can. 1269. \\
\textbf{N11} & \textbf{Pyx} & GIRM 327--329, 333; USCCB GIRM 329; CIC (1983) can. 939; Redemptionis Sacramentum 117--119, 133; Holy Communion and Worship of the Eucharist Outside Mass (1973). 1962 counterpart: Rituale Romanum (1925/1952), Title IV, ch. 2; Missale Romanum (1962), Ritus Servandus X.6. \\
\textbf{N12} & \textbf{Communion Cup} & GIRM 279, 281--287, 327--330, 333; USCCB GIRM 283; USCCB Norms for Holy Communion under Both Kinds (2011) 32--38, 52--55; Redemptionis Sacramentum 100--107, 117--119. 1962 counterpart: Missale Romanum (1962), Rubricae Generales 503; Ritus Servandus X.6. Symbolism: GIRM 281; USCCB Norms for Holy Communion under Both Kinds (2011) 32--33; Matt 26:27--28; 1 Cor 10:16; 1 Cor 11:25--26. \\
\textbf{N13} & \textbf{Cruets} & GIRM 118c, 139--142, 163, 178, 190, 279, 306; CIC (1983) can. 924 §3. 1962 counterpart: Missale Romanum (1962), Rubricae Generales 528; Ritus Servandus VII.4, VII.6, VII.10. \\
\textbf{N14} & \textbf{Flagon} & GIRM 73, 118c, 140, 306, 328--330; Redemptionis Sacramentum 106, 117; USCCB Norms for Holy Communion under Both Kinds (2011) 35--38. 1962 counterpart: Missale Romanum (1962), Rubricae Generales 528 (negative witness: small wine and water cruets only). \\
\textbf{N15} & \textbf{Altar Cloth} & GIRM 117, 304 (with USCCB territorial adaptation); GIRM 280, 306; Redemptionis Sacramentum 120. 1962 counterpart: Missale Romanum (1962), Rubricae Generales 526. Symbolism: GIRM 304; St. Thomas Aquinas, Summa Theologiae III, q. 83, a. 3, ad 7; Matt 27:59. \\
\textbf{N16} & \textbf{Corporal} & GIRM 118c, 139, 163, 178, 183, 190, 192, 279--280, 306; Redemptionis Sacramentum 120; USCCB Norms for Holy Communion under Both Kinds (2011) 53. 1962 counterpart: Missale Romanum (1962), Ritus Servandus I.1, II.2, VI.7. Symbolism: St. Thomas Aquinas, Summa Theologiae III, q. 83, a. 3, ad 7 (citing De Consecr. dist. 1); Matt 27:59; John 20:6. \\
\textbf{N17} & \textbf{Purificator} & GIRM 118c, 139, 163, 178, 192, 279, 286, 306; Redemptionis Sacramentum 119--120; USCCB Norms for Holy Communion under Both Kinds (2011) 43, 53. 1962 counterpart: Missale Romanum (1962), Ritus Servandus I.1, VII.4, X.5. \\
\textbf{N18} & \textbf{Pall} & GIRM 118c (pro opportunitate), 139, 306; Redemptionis Sacramentum 120. 1962 counterpart: Missale Romanum (1962), Ritus Servandus I.1, VII.5, VIII.7. \\
\textbf{N19} & \textbf{Lavabo Towel} & GIRM 76, 118c, 145, 190. 1962 counterpart: Missale Romanum (1962), Rubricae Generales 528; Ritus Servandus VII.6--7, VII.10; Ordo Missae 1039. \\
\textbf{N20} & \textbf{Roman Missal} & GIRM 30, 118a, 139, 147, 189, 306, 349, 389; Postquam Summus Pontifex (2021) 2--3, 40. 1962 counterpart: Missale Romanum (1962), Rubricae Generales 527; Ritus Servandus I.1, II.4--5, VI.5, XI.1. Symbolism: GIRM 349. \\
\textbf{N21} & \textbf{Lectionary for Mass} & GIRM 28, 118b, 120d, 172, 194, 349; Ordo lectionum Missae (1981) Praenotanda 35--37, 69, 103--107. 1962 counterpart: Missale Romanum (1962), Ritus Servandus VI.4. Symbolism: GIRM 28, 349; Ordo lectionum Missae (1981) Praenotanda 35; Ps 118:105 [119:105]. \\
\textbf{N22} & \textbf{Book of the Gospels} & GIRM 62, 117, 120d, 122, 133--134, 172--175, 194--195, 277, 306, 349; Ordo lectionum Missae (1981) Praenotanda 36--37. 1962 counterpart: Missale Romanum (1962), Ritus Servandus VI.5. Symbolism: GIRM 62, 134, 175, 349; St. Thomas Aquinas, Summa Theologiae III, q. 83, a. 4, corpus. \\
\textbf{N23} & \textbf{Book of the Chair} & GIRM 118a, 306; USCCB Excerpts from the Roman Missal (CDWDS confirmed, 2018). 1962 counterpart: No presbyteral counterpart in Missale Romanum (1962), Rubricae Generales 527; episcopal book-bearer precedent only in Caeremoniale Episcoporum (1948) I.20. \\
\textbf{N24} & \textbf{Missal Stand} & GIRM 139, 306--307 (practical custom supporting the altar Missal). 1962 counterpart: Missale Romanum (1962), Rubricae Generales 527; Ritus Servandus VI.5, XI.1. \\
\textbf{N25} & \textbf{Processional Cross} & GIRM 117; GIRM 120b; GIRM 122; GIRM 188; GIRM 274; GIRM 308. 1962 counterpart: Missale Romanum (1962), Rubricae Generales n. 527; Fortescue (1943), pp. 17, 71, 78--79. Symbolism: GIRM 308; 1 Cor 1:23; Gal 6:14; Col 1:20; St. Thomas Aquinas, Summa Theologiae III, q. 83, a. 5, ad 3. \\
\textbf{N26} & \textbf{Altar Candles} & GIRM 117; GIRM 120b; GIRM 122; GIRM 133; GIRM 175; GIRM 274; GIRM 307. 1962 counterpart: Missale Romanum (1962), Rubricae Generales n. 527; Ritus Servandus II.5, VI.5; Fortescue (1943), pp. 17, 89--90. Symbolism: GIRM 117, 307; John 8:12. \\
\textbf{N27} & \textbf{Sanctuary Bells} & GIRM 150. 1962 counterpart: Missale Romanum (1962), Rubricae Generales n. 528; Ritus Servandus VII.8, VIII.6, X.6. \\
\textbf{N28} & \textbf{Aspergillum and Aspersorium} & GIRM 51; GIRM 118c; Roman Missal (2008/2011) Appendix II; Tres abhinc annos (1967) 26. 1962 counterpart: Missale Romanum (1962), Ordo ad aspergendam aquam benedictam, pp. [231]--[233]; Fortescue (1917), p. 16; Fortescue (1943), pp. 79--81. Symbolism: GIRM 51; Roman Missal, Appendix II; Ps 50:9 [Heb 51:7]. \\
\textbf{N29} & \textbf{Thurible} & GIRM 120a; GIRM 133; GIRM 144; GIRM 150; GIRM 175; GIRM 179; GIRM 276; GIRM 277. 1962 counterpart: Missale Romanum (1962), Ritus Servandus IV.4, VI.5, VII.10, VIII.8; Ordo Missae nn. 1035--1038. Symbolism: GIRM 276; Ps 140:2 [Heb 141:2]; Apoc 5:8, 8:3--4; 2 Cor 2:14; St. Thomas Aquinas, Summa Theologiae III, q. 83, a. 5, ad 2. \\
\textbf{N30} & \textbf{Incense Boat} & GIRM 120; GIRM 144; GIRM 276; GIRM 277; Caeremoniale Episcoporum (1984). 1962 counterpart: Missale Romanum (1962), Ritus Servandus IV.4, VII.10. \\
\textbf{N31} & \textbf{Alb} & GIRM 119; GIRM 335; GIRM 336; GIRM 339 (USCCB); Redemptionis Sacramentum 121--122, 126; Caeremoniale Episcoporum (1984) 65. 1962 counterpart: Missale Romanum (1962), Rubricae Generales Tit. XIX.1; Ritus Servandus I.2--3, II.1, VI.8; Praeparatio ad Missam (Ad Albam). Symbolism: GIRM 335--336; Gal 3:27; Apoc 7:13--14; traditional Roman sacristy vesting prayer (Ad Albam); Amalarius of Metz, Liber Officialis II.18. \\
\textbf{N32} & \textbf{Cincture} & GIRM 336; GIRM 340; Redemptionis Sacramentum 121--122; Caeremoniale Episcoporum (1984) 65. 1962 counterpart: Missale Romanum (1962), Ritus Servandus I.3; Praeparatio ad Missam (Ad Cingulum). Symbolism: Practical norm: GIRM 336; devotional tradition: Roman sacristy vesting prayer (Ad Cingulum); Luke 12:35; Eph 6:14; Amalarius of Metz, Liber Officialis II.22. \\
\textbf{N33} & \textbf{Amice} & GIRM 336; Redemptionis Sacramentum 121--122; Caeremoniale Episcoporum (1984) 65. 1962 counterpart: Missale Romanum (1962), Ritus Servandus I.3; Praeparatio ad Missam (Ad Amictum). Symbolism: Practical norm: GIRM 336; devotional tradition: Roman sacristy vesting prayer (Ad Amictum); Eph 6:17; Ps 38:2 [Heb 39:1]; Ps 140:3 [Heb 141:3]; Amalarius of Metz, Liber Officialis II.17. \\
\textbf{N34} & \textbf{Stole} & GIRM 119; GIRM 336; GIRM 337; GIRM 338; GIRM 340; GIRM 346 (USCCB 346e, 346h); Redemptionis Sacramentum 123--126. 1962 counterpart: Missale Romanum (1962), Rubricae Generales n. 137; Ritus Servandus I.3--4; Praeparatio ad Missam (Ad Stolam). Symbolism: GIRM 335, 340; traditional Roman sacristy vesting prayer (Ad Stolam); Matt 11:29--30; Amalarius of Metz, Liber Officialis II.20. \\
\textbf{N35} & \textbf{Chasuble} & GIRM 119a; GIRM 209; GIRM 335; GIRM 337; GIRM 344; GIRM 346 (USCCB 346e, 346h); Redemptionis Sacramentum 123--124, 126--127; Tres abhinc annos (1967) 26. 1962 counterpart: Missale Romanum (1962), Rubricae Generales nn. 133--134, 137; Ritus Servandus I.4; Praeparatio ad Missam (Ad Casulam). Symbolism: GIRM 335, 337; Matt 11:30; Col 3:14; traditional Roman sacristy vesting prayer (Ad Casulam); Amalarius of Metz, Liber Officialis II.19. \\
\textbf{N36} & \textbf{Dalmatic} & GIRM 119b, 335--336, 338, 340, 342--346; Redemptionis Sacramentum 125--126; Caeremoniale Episcoporum (1984) 56. 1962 counterpart: Missale Romanum (1962), Rubricae Generales 134, 137; Caeremoniale Episcoporum (1948) II.8 n. 18. Symbolism: GIRM 335, 338; Caeremoniale Episcoporum (1984) 56; Is 61:10; James 1:27; traditional Roman vesting prayer; Amalarius of Metz, Liber Officialis II.21. \\
\textbf{N37} & \textbf{Cassock and Surplice} & GIRM 119c, 336, 339; USCCB GIRM 339; Redemptionis Sacramentum 128; Caeremoniale Episcoporum (1984) 65. 1962 counterpart: Missale Romanum (1962), Ritus Servandus I.2, II.1, VI.8; Caeremoniale Episcoporum (1948) I.11 §5. \\
\textbf{N38} & \textbf{Cope} & GIRM 337, 341, 346; Holy Communion and Worship of the Eucharist Outside Mass (1973) 91--92; Tres abhinc annos (1967) 26; Caeremoniale Episcoporum (1984). 1962 counterpart: Missale Romanum (1962), Rubricae Generales Cap. XIX (137); Caeremoniale Episcoporum (1948) II.8 n. 18. \\
\textbf{N39} & \textbf{Humeral Veil} & Holy Communion and Worship of the Eucharist Outside Mass (1973) 91--92; Roman Missal (2002/2008/2011), Holy Thursday Evening Mass of the Lord's Supper; Caeremoniale Episcoporum (1984); Ministeria quaedam (1972). 1962 counterpart: Missale Romanum (1962), Ritus Servandus VII.9, X.8; Holy Thursday Transfer. Symbolism: Holy Communion and Worship of the Eucharist Outside Mass (1973) nn. 91--92; Roman Missal (2008/2011), Evening Mass of the Lord's Supper. \\
\textbf{N40} & \textbf{Holy Oils and Oil Stocks} & CIC 1983 can. 847, pars. 1--2; CIC 1983 can. 999 n. 2; can. 1003, par. 1; Ordo benedicendi oleum catechumenorum et infirmorum et conficiendi chrisma (3 Dec 1970; editio typica 1971); Caeremoniale Episcoporum (14 Sept 1984) no. 274; Missale Romanum (2002/2008; US 2011), Missa Chrismatis; CCC nos. 1241, 1294--1300. 1962 counterpart: Missale Romanum (1962), Feria V in Coena Domini (Missa Chrismatis) and Vigilia Paschalis (Blessing of Baptismal Water); Pontificale Romanum (1859), De benedictione oleorum. Symbolism: CCC 1294--1300; Ordo benedicendi oleum (1970); Roman Chrism Preface; 2 Cor 1:21--22; James 5:14--15. \\
\textbf{N41} & \textbf{Subdeacon's Tunicle and Maniple} & Tres abhinc annos (4 May 1967, AAS 59:442--448; eff. 29 June 1967) no. 25; Missale Romanum, editio typica prima (26 Mar 1970, Prot. 166/70), IGMR nos. 142--152, 301; Pope Paul VI, Ministeria quaedam (15 Aug 1972, AAS 64:529--534; eff. 1 Jan 1973) Norms I--IV; Missale Romanum, editio typica altera (27 Mar 1975); Caeremoniale Episcoporum (14 Sept 1984) no. 56; GIRM (2002/2008; US 2011) nos. 119, 335--341. 1962 counterpart: Missale Romanum (1962), Rubricae Generales nn. 134, 137; Ritus Servandus I.3--4; Praeparatio ad Missam p. LXXIV; Caeremoniale Episcoporum (1948) II.8 n. 18. Symbolism: Traditional Roman Missal vesting prayers (Praeparatio ad Missam); Ps 125:5--7; Amalarius of Metz, Liber Officialis II.22, II.24. \\
\textbf{N42} & \textbf{Fragile Vessels, Flagon Pouring, and EMHC Purification} & IGMR (1970/1975) nos. 290--291 (historical baseline); GIRM (2002/2008; US 2011) nos. 279, 327--330 (with USCCB adaptation at no. 329); Redemptionis Sacramentum (25 Mar 2004, AAS 96:549--601) nos. 106, 117, 119; CDWDS letter declining renewal of U.S. EMHC purification indult (Oct 2006); USCCB Norms for the Distribution and Reception of Holy Communion under Both Kinds (2002; revised 2011) nos. 35--38, 52--53. 1962 counterpart: Missale Romanum (1962), Ritus Servandus I.1, X.5--6; Rubricae Generales n. 503. \\
\textbf{N43} & \textbf{Sanctuary, Cross, Tabernacle, Altar Card, and Veil Clarifications (1970/1975)} & IGMR (1975) nos. 262, 270, 276 (historical baseline); CDWDS Response on GIRM no. 299 (25 Sept 2000, Prot. N. 2036/00/L; Notitiae 36 [2000] 171--173); GIRM (2002/2008; US 2011) nos. 118c, 299, 306, 308, 314--315; Missale Romanum (1962), Rubricae Generales n. 527; Ritus Servandus I.1. 1962 counterpart: Missale Romanum (1962), Rubricae Generales nn. 526--527; Ritus Servandus I.1, II.2; CIC 1917 cann. 1268--1269. \\
\bottomrule
\end{longtable}}
